\section{The Claim}

The full derivation of the base belongs to the theory paper \cite{pollard2026vibe}. What the rest of this paper rests on is stated here once and used throughout.

Vibe is a discrete, deterministic, reversible substrate that is conjectured to coarse-grain toward a conserved-current field theory of the kind the continuum programs describe, with no fitted continuum constant and no randomness. The same structures those programs reach from the top, vibe builds from the bottom. This is a conjecture with one firm anchor, not a finished derivation.

The model rests on two axioms. The first names the substance, the second names its motion.

\begin{axiom}[The vibe]
The one substance is the vibe, which is experience. A vibe's tone is the quality of that experience, not a thing apart from it. There is nothing beneath it.
\end{axiom}

\begin{axiom}[The flow]
The whole is the flow, the eight base elements running together as one, the vibe in endless motion.
\end{axiom}

The first axiom is the one place the model reaches past physics. Nothing downstream depends on the substance being felt. A reader who wants only the physics may read the vibe as a plain three-valued charge and lose nothing. We state it because the structure that yields the physics is, on the fuller reading, the structure of experience itself.

The base is exactly eight elements, no more. Three give the where, two the what, two the how, and one the when. Each is a four-letter word, so each doubles as a clean symbol.

\begin{table}[ht]
\centering
\small
\begin{tabularx}{\textwidth}{@{}llX@{}}
\toprule
element & role & what it is \\
\midrule
\texttt{mesh} & where & the whole space, the growing $\{3,4,3,4\}$ hyperbolic honeycomb of 24-cells \\
\texttt{dock} & where & one cell, a here-now where tones meet and move on, the 24-cell with twenty-four neighbors \\
\texttt{site} & where & one of a dock's twenty-four directions, a $D_4$ root, holding one tone \\
\texttt{tone} & what & the value at a site, one of $\{-1, 0, +1\}$ \\
\texttt{lean} & what & the order $-1 < 0 < +1$, turning tone into signed charge and giving growth its direction \\
\texttt{knit} & how & the law of the state, collide then stream, once a beat, reversible and charge-conserving \\
\texttt{wake} & how & the law of the space, new docks born quiet at the growing edge, beat after beat \\
\texttt{beat} & when & one discrete tick of time \\
\bottomrule
\end{tabularx}
\caption{The eight base elements. The whole model is these and nothing else.}
\label{tab:elements}
\end{table}

Nothing is ever added underneath the eight. The base is discrete, deterministic, and carries no randomness. A ternary value, a finite direction set, a discrete beat, a rule that is a permutation, and a finite symmetry group. Real numbers appear only as measured outputs, never at the base. The continuum, Lie groups, the Dirac and Maxwell equations, Lorentz invariance, and the smooth metric are all emergent by coarse-graining. Vibes pass through one to three middle layers before they are physics, so a raw tone is not a particle and a first cluster is not a force.

This paper grades every claim on one scale, kept fixed throughout. \textbf{F}, firm, a structure-identical bridge established on the substrate. \textbf{P}, parallel, the same shape reached by an independent route but not yet derived from the base. \textbf{O}, open, a well-posed frontier with a stated falsifying test, unsolved on both sides. The result labels are narrower. \textbf{Derived} is forced from the base with no measured number fed in. \textbf{Parallel} is the same destination by a different road. \textbf{Open} is unsolved with a named obstruction. A separate axis runs alongside. Tier A is derived from the base, Tier B is checked against external numbers, and a Standard Model constant fed in for comparison never counts as a derivation. The two axes are read together and the substrate status governs.
